% =============================================================================
\section{Analysis and Discussion}
\label{sec:analysis}
% =============================================================================

We now look beyond aggregate NDCG@10 numbers and ask \emph{why} the directional pattern in \S\ref{sec:results} arises. Four lenses inform this: position bias (\S\ref{sec:bias}), query difficulty (\S\ref{sec:difficulty}), ranking agreement (\S\ref{sec:agreement}), and a quality-cost Pareto view that frames the refinement target (\S\ref{sec:pareto}).

\subsection{Position Bias Analysis}
\label{sec:bias}

We investigate a question unique to our framework: do LLMs exhibit different position bias patterns when asked to select the worst versus the best document, and---more importantly---does the same model behave differently when ``best'' and ``worst'' are asked \emph{jointly} versus alone?

\textbf{Experimental design}: For each comparison, we record which passage label (A, B, C, or D) was selected, along with the comparison type (\texttt{best}, \texttt{worst}, \texttt{dual\_best}, \texttt{dual\_worst}). We analyze approximately 170K comparisons per model across all methods.

\textbf{Best-selection bias}: For $f_{\text{best}}$ (\topdown{}), we observe a strong U-shaped bias across all models: positions A (first) and D (last) are consistently overselected, while positions B and C are underselected. For Qwen3-14B on DL 2019 (representative): $A=35.1\%$, $B=12.5\%$, $C=11.2\%$, $D=41.1\%$ (expected: 25\% each, $\chi^2{=}5703$, $p \ll 0.001$). The recency bias (D) is the dominant effect, strongest on Flan-T5-XL ($D=53.3\%$) and Qwen3.5-4B ($D=54.7\%$).

\textbf{Worst-selection bias}: For $f_{\text{worst}}$ (\bottomup{}), position D is overwhelmingly selected as ``worst'': D-frequencies range from 36.5\% (Qwen3-4B) to 63.1\% (Flan-T5-XL), far exceeding the 25\% baseline. Positions B and C are severely underselected ($<$17\% each). The pattern is consistent across all 9 models, with $\chi^2$ values exceeding 5000 in every case. This extreme recency bias for worst-selection is the most parsimonious explanation for \bottomup{}'s poor effectiveness: the model defaults to labeling the last passage as ``worst'' regardless of content.

\textbf{Dual-end bias}: The dual prompt produces a \textbf{novel asymmetric pattern}. For \texttt{dual\_best}, the bias is substantially \emph{reduced} compared to standalone best-selection---Qwen3-14B shows $A=22.2\%$, $B=26.8\%$, $C=18.0\%$, $D=33.1\%$, much more uniform ($\chi^2{=}1683$ vs.\ $5703$ for standalone best). For \texttt{dual\_worst}, the pattern \emph{reverses}: position A (first) becomes the most frequently selected as worst (41.3\%), while D (last) is the \emph{least} selected (12.4\%). This reversal is consistent across all 9 models on both DL19 and DL20.

\begin{table}[!htbp]
\caption{Position D (last passage) selection frequency by comparison type, averaged over 9 models on DL19. The middle column shows that \bottomup{} (\texttt{worst}) collapses heavily onto position D, while \dualend{} (\texttt{dual\_worst}) does the opposite.}
\label{tab:posD}
\small
\begin{tabular}{lcccc}
\toprule
Type & A & B & C & D \\
\midrule
\texttt{best} & .33 & .13 & .13 & .41 \\
\texttt{worst} & .24 & .10 & .13 & .50 \\
\texttt{dual\_best} & .27 & .24 & .16 & .33 \\
\texttt{dual\_worst} & .35 & .22 & .28 & .13 \\
\bottomrule
\end{tabular}
\end{table}

\textbf{Interpretation}: The reversed \texttt{dual\_worst} pattern suggests that when asked for both best \emph{and} worst simultaneously, the model behaves differently than when asked for either signal alone. One natural reading is two-stage: the model first identifies the best (still benefiting from a recency tendency toward D), then assigns ``worst'' to the most salient remaining candidate---which is now typically the first passage (primacy effect). Either way, the empirical finding is that joint elicitation is not equivalent to running standalone best and standalone worst as parallel queries; the prompt structure changes the answer distribution.

\textbf{Implications}: The dual prompt acts as an implicit debiasing mechanism. Unlike permutation voting~\citep{tang2024found}, which requires multiple passes to reduce bias, dual-end achieves bias reduction \emph{within a single prompt}. We see this as the strongest mechanistic clue in the paper, and it directly motivates the order-robust ``bias-aware DualEnd'' refinement variant (\S\ref{sec:refinement}), which exploits the joint-prompt asymmetry by running a small number of controlled orderings only on hard windows.


\subsection{Query Difficulty Stratification}
\label{sec:difficulty}

We stratify queries by difficulty (measured by BM25 NDCG@10) into Easy / Medium / Hard terciles and report $\Delta$NDCG@10 of \bottomup{}-Heap and \dualend{}-Cocktail against \topdown{}-Heap. Table~\ref{tab:difficulty} shows the pattern for Qwen3-14B on DL 2019, which is representative of the larger Qwen models.

\begin{table}[!htbp]
\caption{$\Delta$NDCG@10 vs.\ \topdown{}-Heap, stratified by BM25 difficulty terciles (Qwen3-14B, DL 2019, 43 queries). \dualend{}-Cocktail helps most on Medium and Hard queries; \bottomup{}-Heap hurts most on Hard queries.}
\label{tab:difficulty}
\small
\begin{tabular}{lcccc}
\toprule
Method & Easy & Medium & Hard & Overall \\
\midrule
\bottomup{}-Heap & $-$.030 & $-$.014 & $-$.103 & $-$.048 \\
\dualend{}-Cocktail & $-$.006 & +.019 & +.009 & +.007 \\
\bottomrule
\end{tabular}
\end{table}

For Qwen3-14B, \dualend{}-Cocktail gains the most on medium queries (+.019) and is roughly neutral on easy and hard queries. \bottomup{}-Heap is consistently negative, with the largest hurt on hard queries ($-.103$). For Flan-T5-XL on DL 2019, the pattern is somewhat reversed at the extremes: \dualend{}-Cocktail helps on easy queries (+.014) and is mildly negative on medium and hard (--.010 / --.010), so the difficulty story should be presented as a trend conditional on model rather than as a universal rule. Across all 18 model--dataset configurations, the average DualEnd help is largest on the Easy and Medium terciles and shrinks on Hard, while BottomUp hurts roughly uniformly across terciles. Together these patterns motivate routing rather than universal use---if the gain is concentrated in particular query slices, the right place to spend DualEnd's extra calls is on those slices.


\subsection{Ranking Agreement}
\label{sec:agreement}

To check whether \dualend{} is essentially a perturbation of \topdown{} or a genuinely different ranking, we compute pairwise ranking agreement (top-10 overlap and Kendall's $\tau$) between method pairs. Averaging over all 18 model--dataset configurations:
\begin{itemize}[leftmargin=*]
\item \topdown{} vs.\ \dualend{}: overlap@10 $\approx 7.0$, $\tau \approx 0.93$ (high agreement).
\item \topdown{} vs.\ \bottomup{}: overlap@10 $\approx 5.0$, $\tau \approx 0.86$ (lower agreement).
\item \bottomup{} vs.\ \dualend{}: overlap@10 $\approx 5.2$, $\tau \approx 0.88$.
\end{itemize}

Two facts follow. First, \dualend{} stays mostly aligned with \topdown{}: it preserves the useful TopDown signal and adds refinement from the worst-selection side, rather than producing a structurally different ranking. Second, \bottomup{} is not random---its agreement with \topdown{} is still ``high'' on the standard tau scale---but it is consistently noisier and more divergent than \dualend{}. This is exactly the regime in which fusion is unhelpful: the diverging signal does not encode complementary truth, so adding it to \topdown{} only adds variance, which is what we see in \S\ref{sec:rq3}.


\subsection{When Does DualEnd Help?}
\label{sec:when_helps}

To make the ``modest mean, model-dependent variance'' picture concrete, we ran a per-query help/hurt summary on three representative DL19 model settings using a small qualitative analysis script that operates on the saved run files. Table~\ref{tab:when_helps} reports the mean delta and the help/hurt/tie counts.

\begin{table}[!htbp]
\caption{Per-query help/hurt summary for \dualend{}-Cocktail vs.\ \topdown{}-Heap on three representative DL19 settings (43 queries).}
\label{tab:when_helps}
\small
\begin{tabular}{lrrrr}
\toprule
Model & Mean $\Delta$ & Help & Hurt & Tie \\
\midrule
flan-t5-xl & $-0.0017$ & 20 & 16 & 7 \\
qwen3-8b & $+0.0336$ & 26 & 12 & 5 \\
qwen3-14b & $+0.0072$ & 19 & 15 & 9 \\
\bottomrule
\end{tabular}
\end{table}

Three observations follow. First, DualEnd is not universally better than TopDown: on flan-t5-xl the mean delta is mildly negative even though DualEnd ``wins'' more queries than TopDown (20 vs.\ 16) than ``ties'' or ``loses.'' Second, the strongest positive regime in our setting is Qwen-style causal models, especially the smaller ones. Third, on the strongest setting (qwen3-8b), most of the positive cases correspond to genuine relevant-document additions to the top-$k$, while on flan-t5-xl the ``hurt'' cases lean more toward relevant-document drops or noisy mixed changes. This directly motivates routing the joint signal where it is likely to matter rather than spending it everywhere.


\subsection{Quality--Cost Pareto Frontier}
\label{sec:pareto}

We compute a Pareto analysis over the existing TREC DL run folders so the paper can talk about budgeted quality rather than isolated averages. For each method we average NDCG@10 against three cost axes (comparisons, total tokens, time per query) across all 9 models on DL19.

\begin{table}[!htbp]
\caption{Global mean Pareto frontier members on TREC DL 2019 (NDCG@10 vs.\ each cost axis, averaged across all 9 models). The comparison- and token-frontier contains only three points: \topdown{}-Heap, \topdown{}-Bubble, and \dualend{}-Cocktail.}
\label{tab:pareto}
\small
\begin{tabular}{lc}
\toprule
Frontier & Members \\
\midrule
Comparisons & TD-Heap, TD-Bubble, DE-Cocktail \\
Total tokens & TD-Heap, TD-Bubble, DE-Cocktail \\
Time & PermVote($p$=2), TD-Heap, TD-Bubble, DE-Cocktail \\
\bottomrule
\end{tabular}
\end{table}

\textbf{Frontier members and means} (mean NDCG@10 / mean comparisons / mean total tokens / mean time):
\begin{itemize}[leftmargin=*]
\item \topdown{}-Heap: $.6851$ / $76.5$ / $32464$ / $28.41$ s
\item \topdown{}-Bubble: $.6897$ / $300.0$ / $126456$ / $110.93$ s
\item \dualend{}-Cocktail: $.6962$ / $546.0$ / $233192$ / $212.62$ s
\end{itemize}

\textbf{Interpretation.}
\topdown{}-Heap remains the cheap anchor; \topdown{}-Bubble is the strongest middle-cost anchor; \dualend{}-Cocktail is the quality-first anchor. Crucially, the global frontier contains \emph{no} \bottomup{} or \bidir{} variant on either the comparisons or the tokens axis. \bottomup{} is dominated because it spends many more comparisons than \topdown{}-Heap at strictly worse mean quality; \bidir{} is dominated for the same reason---it inherits BottomUp's noise and cost. The time-axis frontier additionally admits permutation voting ($p{=}2$) as a fast-but-weaker point; the quality-axis story does not change. This identifies the paper's refinement target precisely: the empty region between \topdown{}-Bubble and \dualend{}-Cocktail. The Selective DualEnd, bias-aware DualEnd, and same-call worst-signal regularization variants in \S\ref{sec:refinement} are designed to land in that region rather than to push the upper-right corner further.


\subsection{Error Analysis}
\label{sec:errors}

We close with a brief qualitative error analysis to characterize what each method gets wrong.

\textbf{Top-down errors}: \topdown{} tends to make errors when multiple documents are similarly relevant, and position bias causes the LLM to consistently favor certain positions. This leads to ``promoted impostors''---documents that are ranked higher than they deserve because they were systematically selected as best when placed in favorable positions.

\textbf{Bottom-up errors}: \bottomup{} tends to make errors when documents are similarly \emph{ir}relevant, and the LLM struggles to identify the ``most irrelevant'' among a set of non-relevant documents. Combined with the strong recency bias for the standalone worst-selection prompt, this leads to ``survivor bias''---non-relevant documents that survive elimination and contaminate the top-$k$.

\textbf{Dual-end errors}: \dualend{} introduces a new error type: \emph{interference errors}, where the simultaneous best-and-worst assessment causes the LLM to produce an inconsistent pair (e.g., selecting the same document as both best and worst). Our ``same best\&worst'' warning count shows this occurs in fewer than 0.1\% of dual comparisons (1--2 instances per run), indicating interference is rare. More common are partial parse failures where only one label is extracted ($<$1\% of comparisons), handled by heuristic defaults.

Contrary to our original hypothesis, these complementary error patterns do \emph{not} make \bidir{} fusion effective. The BottomUp errors are too frequent and systematic (driven by position bias, not by random noise) to cancel out TopDown errors. \dualend{}, in contrast, naturally combines both selection directions in a single call, achieving the error-cancellation benefit without the fusion overhead.


\subsection{Relationship to Tournament Theory}
\label{sec:tournament}

Our dual-end selection sort can be viewed through the lens of tournament theory. In a standard single-elimination tournament, each match produces one winner and one loser. The winner advances; the loser is discarded. This is analogous to standard setwise ranking. A \emph{double-elimination tournament} tracks both the winners' bracket and the losers' bracket, giving each competitor two chances before elimination. Our \dualend{}-Selection is similar: each comparison round identifies both bracket winners and bracket losers, and both results are used for ranking. Recent work on tournament graphs for LLM ranking~\citep{blitzrank2026} also seeks to extract richer information from each comparison, but through graph-theoretic analysis rather than prompt-level modifications.

The information-theoretic advantage of dual-end ($\sim$$74$--$90$\% more bits per call depending on window size) provides a theoretical upper bound on the efficiency improvement. In practice, the realized improvement depends on: (1) the accuracy of worst-selection relative to best-selection, (2) the overhead of parsing dual outputs, and (3) whether the sorting algorithm can fully exploit both outputs. Our empirical results in \S\ref{sec:rq2} confirm the practical gain is real but modest---enough to make \dualend{} the strongest family overall, not enough to claim Bonferroni-significant wins on most TREC DL configurations.
